% 2027 Commercial & Investment Bank Risk Management Summer Analyst Program
\documentclass[10pt,letterpaper]{article}

\usepackage[margin=0.75in]{geometry}
\usepackage[hidelinks]{hyperref}
\usepackage[skip=4pt]{parskip}

\pagestyle{empty}
\pagenumbering{gobble}
\setlength{\parindent}{0pt}

\begin{document}

\textbf{Rojae Mighty}\\
Stony Brook, NY \quad$\cdot$\quad
\href{mailto:rojae.mighty@stonybrook.edu}{rojae.mighty@stonybrook.edu}\\
\href{https://rojaemighty.com}{rojaemighty.com} \quad$\cdot$\quad
\href{https://www.linkedin.com/in/rojaemight/}{linkedin.com/in/rojaemight}

\vspace{8pt}
July 14, 2026

\vspace{8pt}
Hiring Manager\\
JPMorganChase\\
New York, NY

\vspace{8pt}
\textbf{Re: 2027 Commercial \& Investment Bank Risk Management Summer Analyst Program}

\vspace{4pt}
Dear Hiring Manager,

I am excited to apply for the 2027 Commercial \& Investment Bank Risk Management Summer Analyst Program at JPMorganChase. As an undergraduate physics student at Stony Brook University, I have developed strong quantitative, analytical, and communication skills through research in nuclear and particle physics, machine learning, and financial data analysis. I am eager to apply this background to analyzing risk, building models, and supporting informed decisions across JPMorganChase's global businesses.

Through my work at Brookhaven National Laboratory, including the Department of Energy Science Undergraduate Laboratory Internship, I studied electron-ion collisions using Monte Carlo event generation, ROOT, Python, and Unix/Linux scripting. I implemented a spatially dependent parton distribution function and analyzed how nuclear environments affect particle behavior. This work required careful data analysis, computational modeling, attention to detail, and the ability to draw meaningful conclusions from complex and uncertain systems---skills directly relevant to due diligence, portfolio monitoring, and risk analysis.

As an undergraduate researcher at Stony Brook's Center for Frontiers in Nuclear Science, I use C++, ROOT, Python, and machine-learning methods to analyze large-scale collision datasets and reconstruct physical quantities. I have also developed a market microstructure forecasting project using Python, pandas, NumPy, scikit-learn, feature engineering, and walk-forward validation on limit order book data. In addition, through J.P. Morgan's Quantitative Research program on Forage, I gained experience with derivatives pricing, statistical data analysis, and credit-risk modeling. Together, these experiences have strengthened my ability to work with data, evaluate assumptions, and communicate technical findings clearly.

I am particularly drawn to this program's emphasis on assessing counterparties and products, monitoring portfolio sensitivities, evaluating market movements, and preparing recommendations for senior management. JPMorganChase's collaborative, global approach to managing credit, market, liquidity, and country risk would provide an exciting environment in which to develop my financial expertise while contributing my perspective as a quantitative problem solver.

Thank you for considering my application. I would welcome the opportunity to discuss how my research experience, technical background, and interest in financial risk management could contribute to the success of JPMorganChase's Commercial \& Investment Bank.

\vspace{4pt}
Sincerely,\\[12pt]
Rojae Mighty

\end{document}
